\documentclass[10pt,a4paper]{amsart}
\usepackage[margin=19mm]{geometry}
\usepackage{amsmath,amssymb,mathtools}
\usepackage[scr=rsfs,cal=euler]{mathalfa}
\usepackage{xfrac}
\usepackage[unicode,hidelinks,bookmarks=false]{hyperref}
\newcommand{\ie}{i.e.\ }
\newcommand{\cf}{cf.\ }
\newcommand{\resp}{respectively\ }
\newcommand{\Subsection}{\subsection}
\providecommand{\nobreakdash}{\nobreak\hbox{-}}
\setlength{\emergencystretch}{2em}
\multlinegap0pt
\usepackage{amssymb}
\usepackage[shortlabels]{enumitem}
\usepackage{enumerate}
\usepackage{mathtools}
\usepackage{bm}
\usepackage{xcolor}
\newtheorem{proposition}{Proposition}
\newtheorem{theorem}{Theorem}
\newcommand\E{\mathbb{E}}
\newcommand\eps{\varepsilon}
\title{Local rigidity of self-joinings and factors of pro-nilsystems}
\author{Pauwel Van Den Eeckhaut, Asgar Jamneshan}
\date{Source: arXiv:2604.02089v2 (CC BY 4.0)}
\begin{document}
\maketitle

\noindent\emph{Source and licence.} Local rigidity of self-joinings and factors of pro-nilsystems. Version arXiv:2604.02089v2 (CC BY 4.0), distributed by arXiv under the Creative Commons Attribution 4.0 International licence, http://creativecommons.org/licenses/by/4.0/, as recorded in the arXiv licence metadata for this exact version and shown on the abstract page https://arxiv.org/abs/2604.02089v2. Full licence notice: \texttt{licenses/arxiv-2604.02089-CC-BY-4.0.txt}.\par\smallskip

\noindent\textit{Editorial scope.} Counted unit: the article's theorem thrm:special\_case (source0.tex lines 913--920). The article's complete original proof of it is delivered with it (source0.tex lines 922--1179, one \texttt{proof} environment of 258 lines). No mathematics is written, repaired or re-derived: the statement and the proof are the article's own lines, in the article's own notation, and no retained line is substituted or re-flowed.  Retained uncounted as the source-local context the proof needs: lines 279--287; lines 626--630; lines 906--909.  Nothing was omitted from any retained span, so the package carries no omission record.  No retained line contains a citation, so the wrapper carries no bibliography and no bibliography entry.  No retained line contains a figure environment, an image-inclusion command or drawing code, so no image is delivered and the package declares no asset dependency.  Editorial formatting only, in addition: an AMS \texttt{amsart} preamble that restates the article's own declarations and macros used by the retained lines, namely the environments \texttt{proposition}, \texttt{theorem} on separate counters, together with the standard packages those lines need.  The retained environments print in this document as Proposition 1, Theorem 1 in order of appearance, whereas the article prints the same items with the numbering of its own full document.  The complete native source of the article as distributed in the arXiv e-print for this exact version is bundled unchanged as \texttt{sources/197709/source0.tex}.
\begin{proposition}\label{prop:neighbourhood_diag_joining_graph}
Let $(X,\mu,T)$ be an ergodic nilsystem. Let $\mu_\Delta := (\mathrm{Id}_X,\mathrm{Id}_X)_*\mu$ be the diagonal joining. Then there exists $\delta>0$ such that, for every
$\lambda\in J_e(X,X)$, if\footnote{We view $J_e(X,X)$ as a subspace of the space $\mathrm{Pr}(X\times X)$  of Borel probability measures on the compact Hausdorff space $X\times X$ which, by the Riesz representation theorem, is naturally identified with a weak-* compact subset of the dual of $C(X\times X)$. 
Since $X\times X$ is metrizable, this topology is metrizable on $\mathrm{Pr}(X\times X)$, and we let $d$ be any metric inducing it.} $d(\lambda,\mu_\Delta)<\delta$, then $\lambda$ is the graph joining of an automorphism\footnote{By an automorphism of $(X,\mu,T)$ we mean an automorphism in the measurable sense, that is, an invertible measurable factor map from $(X,\mu,T)$ to itself, defined up to almost-everywhere equality.} of $(X,\mu,T)$. 

Equivalently, every ergodic self-joining of $(X,\mu,T)$ which is not the
graph joining of an automorphism lies at distance at least $\delta$ from the
diagonal joining. 
\end{proposition}

\begin{proposition}\label{prop:joining_nilsystems_is_nilsystem}
    Let $(X = G/\Gamma, \mu, T)$ and $(X' = G'/\Gamma', \mu', T')$ be $k$-step nilsystems, and let $\lambda$ be an ergodic joining of these systems.
    Then there exists a $k$-step subnilsystem $(Y,T\times T')$ of $(X\times X',T\times T')$ such that $\lambda$ is the Haar measure on $Y$.
    In particular, the system defined by the joining is itself a $k$-step nilsystem.
\end{proposition}

\begin{theorem}\label{thrm:factor_of_nilsystem_is_nilsystem}
    Let $k \geq 1$.
    Then every factor of an ergodic $k$-step nilsystem is again an ergodic $k$-step nilsystem.
\end{theorem}

\begin{theorem}[Compact abelian group extension case]\label{thrm:special_case}
    Let $(X,\mu,T)$ be an ergodic $k$-step pro-nilsystem, and let
    \[
    \pi \colon (X,\mu,T)\to (Y,\nu,S)
    \]
    be a factor map such that $(X,\mu,T)$ is a skew-product extension of $(Y,\nu,S)$ by a compact abelian group $K$.
    Then $(Y,\nu,S)$ is a $k$-step pro-nilsystem.
\end{theorem}

\begin{proof}
    Write
    \[
    (X,\mu,T)=\varprojlim (X_i,\mu_i,T_i),
    \]
    where each $(X_i,\mu_i,T_i)$ is an ergodic $k$-step nilsystem, and let
    \[
    p_i\colon X\to X_i
    \]
    denote the factor maps.

    By assumption, $X$ is of the form 
    \[
    (X,\mu,T)=(Y\times K,\nu\times m_K,T_\rho),
    \]
    where $m_K$ is Haar measure on $K$, $\rho\colon Y\to K$ is a measurable map, and
    \[
    T_\rho(y,g)=(Sy,\rho(y)+g).
    \]
    For $u\in K$, write
    \[
    V_u\colon X\to X,\qquad V_u(y,g)=(y,u+g),
    \]
    for the corresponding vertical rotation.

    For each $i$ and each $u\in K$, define a measure on $X_i\times X_i$ by
    \[
    \lambda_u^{(i)}:=(p_i,p_i\circ V_u)_*\mu.
    \]
    Since $p_i$ is a factor map, $\lambda_u^{(i)}$ is a self-joining of $(X_i,\mu_i,T_i)$.
    Moreover, the system
    \[
    (X_i\times X_i,\lambda_u^{(i)},T_i\times T_i)
    \]
    is a factor of the ergodic system $(X,\mu,T)$ via the map $(p_i,p_i\circ V_u)$, and is therefore ergodic.
    Thus
    \[
    \lambda_u^{(i)}\in J_e(X_i,X_i).
    \]

    For each $i$, define
    \[
    H_i:=\{u\in K : \lambda_u^{(i)} \text{ is the graph joining of an automorphism of }X_i\}.
    \]
    We claim that $H_i$ is an open subgroup of $K$.

    Let $u,v\in H_i$.
    Then there exist automorphisms $\phi_u,\phi_v\colon X_i\to X_i$ such that
    \[
    p_i\circ V_u=\phi_u\circ p_i \qquad \mu\text{-a.e.}
    \]
    and
    \[
    p_i\circ V_v=\phi_v\circ p_i \qquad \mu\text{-a.e.}
    \]
    Hence
    \[
    p_i\circ V_{u+v}
    %=p_i\circ V_u\circ V_v
    %=\phi_u\circ p_i\circ V_v
    =\phi_u\circ \phi_v\circ p_i
    \qquad \mu\text{-a.e.}
    \]
    so $\lambda_{u+v}^{(i)}$ is the graph joining associated to the automorphism $\phi_u\circ\phi_v$.
    Thus $u+v\in H_i$.

    Similarly,
    \[
    \phi_u^{-1}\circ p_i
    =\phi_u^{-1}\circ p_i\circ V_u\circ V_{-u}
    =\phi_u^{-1}\circ \phi_u\circ p_i\circ V_{-u}
    =p_i\circ V_{-u}
    \qquad \mu\text{-a.e.},
    \]
    and therefore $\lambda_{-u}^{(i)}$ is the graph joining associated to $\phi_u^{-1}$.
    Hence $-u\in H_i$, and so $H_i$ is a subgroup of $K$.

    By Proposition~\ref{prop:neighbourhood_diag_joining_graph}, there exists an open neighbourhood
    \[
    U_i\subseteq J_e(X_i,X_i)
    \]
    of the diagonal joining such that every element of $U_i$ is the graph joining of an automorphism of $X_i$.
    
    We claim that the map
\[
K \to J_e(X_i,X_i), \qquad u \mapsto \lambda_u^{(i)}
\]
is continuous.
Let $u_n \to u$ in $K$, and let $f,g \in C(X_i)$.
Then
\[
\int_{X_i\times X_i} f(x)g(x')\, d\lambda^{(i)}_{u_n}
=
\int_X f(p_i(x))\, g(p_i(V_{u_n}x))\, d\mu(x).
\]
Write $F = f\circ p_i$ and $G = g\circ p_i$, viewed as elements of $L^2(X,\mu)$.
Since the vertical rotations $V_u$ define a measure-preserving action of the compact group $K$ on $X$, the associated Koopman representation
\[
U_u h := h\circ V_u
\]
on $L^2(X,\mu)$ is strongly continuous.
Hence
\[
\|U_{u_n}G - U_uG\|_{L^2(\mu)} \to 0,
\]
and therefore
\[
\int_X F \cdot U_{u_n}G \, d\mu \to \int_X F \cdot U_uG \, d\mu.
\]
That is,
\[
\int_{X_i\times X_i} f(x)g(x')\, d\lambda^{(i)}_{u_n}
\to
\int_{X_i\times X_i} f(x)g(x')\, d\lambda^{(i)}_u.
\]
Since finite linear combinations of functions of the form $(x,x')\mapsto f(x)g(x')$ are dense in $C(X_i\times X_i)$, it follows that
\[
\lambda^{(i)}_{u_n} \to \lambda^{(i)}_u
\]
in the weak-* topology.

    Since
    \[
    \lambda_0^{(i)}=(p_i,p_i)_*\mu
    \]
    is the diagonal joining, the preimage of $U_i$ under this map is an open neighbourhood of $0$ contained in $H_i$.
    Therefore $H_i$ is an open subgroup of $K$.

    Since $K$ is compact and $H_i$ is open, the quotient $K/H_i$ is finite.
    Choose a finite set $R_i\subseteq K$ of coset representatives such that
    \[
    R_i+H_i=K.
    \]
    By enlarging the sets $R_i$ if necessary, we may assume that $0\in R_i$ for every $i$, and that
    \[
    R_i\subseteq R_j \qquad \text{whenever } i\leq j.
    \]

    For each $u\in H_i$, let
    \[
    V_u^{(i)}\colon (X_i,\mu_i,T_i)\to (X_i,\mu_i,T_i)
    \]
    be an automorphism whose graph joining is $\lambda_u^{(i)}$.
    Then
    \[
    V_u^{(i)}\circ p_i=p_i\circ V_u \qquad \mu\text{-a.e.}
    \]
    and hence
    \[
    V_u^{-1}(p_i^{-1}(\mathcal X_i)) =_\mu p_i^{-1}(\mathcal X_i),
    \]
    where $\mathcal X_i$ denotes the sigma-algebra of the factor $X_i$.
    Thus the sigma-algebra corresponding to $X_i$ is invariant under vertical rotations by elements of $H_i$.

    We now enlarge $X_i$ to a factor invariant under all vertical rotations.
    For each $i$, set
    \[
    \widetilde X_i:=X_i^{R_i}=\prod_{r\in R_i} X_i,
    \]
    and define
    \[
    q_i\colon X\to \widetilde X_i,\qquad
    q_i(x):=(p_i(V_r x))_{r\in R_i}.
    \]
    Let
    \[
    \widetilde\mu_i:=(q_i)_*\mu,
    \qquad
    \widetilde T_i((x_r)_{r\in R_i}):=(T_i x_r)_{r\in R_i}.
    \]
    Then
    \[
    q_i\colon (X,\mu,T)\to (\widetilde X_i,\widetilde\mu_i,\widetilde T_i)
    \]
    is a factor map, so $(\widetilde X_i,\widetilde\mu_i,\widetilde T_i)$ is ergodic.
    Moreover, it is an ergodic finite self-joining of the nilsystem $X_i$, and hence is itself a $k$-step nilsystem by repeated application of Proposition~\ref{prop:joining_nilsystems_is_nilsystem}.

    The sigma-algebra corresponding to $\widetilde X_i$ is
    \[
    q_i^{-1}(\widetilde{\mathcal X}_i)
    = \bigvee_{r\in R_i} V_r^{-1}(p_i^{-1}(\mathcal X_i)).
    \]
    We claim that this sigma-algebra is invariant under all vertical rotations.
    Let $u\in K$.
    For each $r\in R_i$, since $R_i+H_i=K$, there exist $s(r)\in R_i$ and $h(r)\in H_i$ such that
    \[
    r+u=s(r)+h(r).
    \]
    Therefore
    \[
    V_u^{-1}(V_r^{-1}(p_i^{-1}(\mathcal X_i)))
    =V_{r+u}^{-1}(p_i^{-1}(\mathcal X_i))
    =V_{s(r)}^{-1}\bigl(V_{h(r)}^{-1}(p_i^{-1}(\mathcal X_i))\bigr)
    =_\mu V_{s(r)}^{-1}(p_i^{-1}(\mathcal X_i)),
    \]
    using the $H_i$-invariance of $p_i^{-1}(\mathcal X_i)$.
    Hence
    \[
    V_u^{-1}(q_i^{-1}(\widetilde{\mathcal X}_i))
    =_\mu \bigvee_{r\in R_i} V_{s(r)}^{-1}(p_i^{-1}(\mathcal X_i))
    \subseteq_\mu q_i^{-1}(\widetilde{\mathcal X}_i).
    \]
    Since $V_u$ is invertible, the reverse inclusion also holds, and thus
    \[
    V_u^{-1}(q_i^{-1}(\widetilde{\mathcal X}_i)) =_\mu q_i^{-1}(\widetilde{\mathcal X}_i)
    \]
    for every $u\in K$.

    Since $0\in R_i$, each factor $\widetilde X_i$ lies above $X_i$.
    Moreover, because the families $R_i$ and $X_i$ are increasing, the sigma-algebras $q_i^{-1}(\widetilde{\mathcal X}_i)$ are increasing.

    For each $i$, define
    \[
    \mathcal W_i
    := \overline{\pi^{-1}(\mathcal Y)}^{\,\mu}\cap
       \overline{q_i^{-1}(\widetilde{\mathcal X}_i)}^{\,\mu},
    \]
    where the bar denotes $\mu$-completion.
    Let $W_i$ be the corresponding factor of $X$.
    By construction, $W_i$ is a factor of $\widetilde X_i$, and hence, by Theorem~\ref{thrm:factor_of_nilsystem_is_nilsystem}, each $W_i$ is an ergodic $k$-step nilsystem.

    We now show that the sigma-algebra $\pi^{-1}(\mathcal Y)$ is generated by the increasing family $(\mathcal W_i)$.
    Let $f\in L^2(\pi^{-1}(\mathcal Y))$ and $\eps>0$.     Since $X=\varprojlim X_i$ and $\widetilde X_i$ lies above $X_i$, there exist $i_0$ and $ g\in L^2(q_{i_0}^{-1}(\widetilde{\mathcal X}_{i_0}))$ 
    such that $ \|g-f\|_{L^2(\mu)}<\eps$. 
    Set
    \[
    h:=\E_\mu(g\mid \pi^{-1}(\mathcal Y)).
    \]
    Since $X$ is an extension of $Y$ by $K$, we have
    \[
    h(x)=\int_K g(V_u x)\,dm_K(u)
    \qquad \mu\text{-a.e.}
    \]
    Because $q_{i_0}^{-1}(\widetilde{\mathcal X}_{i_0})$ is invariant under all vertical rotations, each function $g\circ V_u$ is measurable with respect to this sigma-algebra.
    Hence $h$ is also measurable with respect to $q_{i_0}^{-1}(\widetilde{\mathcal X}_{i_0})$.
    By construction, $h$ is $\pi^{-1}(\mathcal Y)$-measurable as well.
    Therefore $h\in L^2(\mathcal W_{i_0})$. 

    Moreover,
    \[
    \|h-f\|_{L^2(\mu)}
    =\|\E_\mu(g-f\mid \pi^{-1}(\mathcal Y))\|_{L^2(\mu)}
    \leq \|g-f\|_{L^2(\mu)}
    <\eps.
    \]
    Since $f$ and $\eps$ were arbitrary, it follows that
    \[
    \pi^{-1}(\mathcal Y)\subseteq_\mu \bigvee_i \mathcal W_i.
    \]
    The reverse inclusion is immediate from the definition of $\mathcal W_i$.
    Hence
    \[
    \pi^{-1}(\mathcal Y)=_\mu \bigvee_i \mathcal W_i.
    \]

    Thus the factor $(Y,\nu,S)$ is the inverse limit, in the measure-theoretic sense, of the increasing family of $k$-step nilsystems $W_i$.
    Therefore $(Y,\nu,S)$ is a $k$-step pro-nilsystem.
\end{proof}

\end{document}
